\documentclass[10pt,a4paper]{amsart}
\usepackage[margin=19mm]{geometry}
\usepackage{amsmath,amssymb,mathtools}
\usepackage[scr=rsfs,cal=euler]{mathalfa}
\usepackage{xfrac}
\usepackage[unicode,hidelinks,bookmarks=false]{hyperref}
\newcommand{\ie}{i.e.\ }
\newcommand{\cf}{cf.\ }
\newcommand{\resp}{respectively\ }
\newcommand{\Subsection}{\subsection}
\providecommand{\nobreakdash}{\nobreak\hbox{-}}
\setlength{\emergencystretch}{2em}
\multlinegap0pt
\usepackage{bm}
\usepackage{bbm}
\usepackage{dsfont}
\usepackage{xspace}
\usepackage{cancel}
\usepackage{mathtools}
\usepackage{cleveref}
\usepackage{amsthm}
\makeatletter
\@ifundefined{theorem}{\newtheorem{theorem}{Theorem}}{}
\@ifundefined{lemma}{\newtheorem{lemma}[theorem]{Lemma}}{}
\@ifundefined{proposition}{\newtheorem{proposition}[theorem]{Proposition}}{}
\makeatother
\DeclareMathOperator{\SL}{SL}
\DeclareMathOperator{\ddiv}{div}
\DeclareMathOperator{\gon}{gon}
\def\torz#1{\mathbb Z /#1 \mathbb Z}
\newcommand{\C}{\mathbb{C}}
\newcommand{\Q}{\mathbb{Q}}
\newcommand{\Z}{\mathbb{Z}}
\newcommand{\diamondop}[1]{\langle #1 \rangle} % diamond operator
\newcommand{\set}[1]{\left\lbrace #1 \right\rbrace}
\title[Classification of torsion of elliptic curves over quartic fi]{Classification of torsion of elliptic curves over quartic fields (Proposition 4)}
\author{Maarten Derickx, Filip Najman}
\date{Source: arXiv:2412.16016v2 (CC BY 4.0)}
\begin{document}
\maketitle

\noindent\emph{Source and licence.} Classification of torsion of elliptic curves over quartic fields. Version arXiv:2412.16016v2 (CC BY 4.0), distributed under the Creative Commons Attribution 4.0 International licence, as recorded in the arXiv licence metadata for this exact version and shown on the abstract page https://arxiv.org/abs/2412.16016v2. Full rights notice: licenses/arxiv-2412.16016-CC-BY-4.0.txt.\par\smallskip

\noindent\textit{Editorial scope.} Counted unit: the Proposition printed as Proposition 4 of the article (source0.tex lines 466--483, source label \texttt{prop:hecke\_gon}), delivered together with the article's own complete original proof of it (source0.tex lines 484--511, one \texttt{proof} environment of 28 lines). No mathematics is written, repaired or re-derived: statement and proof are the article's own text line for line. Required context retained uncounted: prop:hecke-kernel (lines 437--442); lem:special\_cusp (lines 448--450). Every definition, display and label of those lines is retained unchanged. Omitted from this package and not redistributed: 1--436, 443--447, 451--465, 512--1412. Nothing inside a retained excerpt is omitted or altered: every retained body is delivered exactly as the article's lines are written, so no omission receipt of any kind accompanies this package. Citations: the retained excerpts carry no citation command at all, so no bibliography entry of the article is transcribed and no reference is redistributed. Reference adaptation: none was needed and none was made; every reference inside the delivered unit resolves inside it, so no unresolved reference or citation is produced. Printed slips kept literal: no printed word, symbol, hypothesis or number of the counted statement or of its proof was altered. Layout and class: the article's own class is \texttt{amsart} with packages including comment, amssymb, tikz-cd, tikz, amsthm, graphicx, caption, subcaption, hyperref, cleveref, booktabs, float, todonotes; the wrapper is the lane's pinned \texttt{amsart} editorial wrapper at 10pt on A4, which restates the theorem-like environments the retained text needs and adds the article's own macro definitions that the retained text needs (\texttt{\textbackslash C}, \texttt{\textbackslash Q}, \texttt{\textbackslash SL}, \texttt{\textbackslash Z}, \texttt{\textbackslash ddiv}, \texttt{\textbackslash diamondop}, \texttt{\textbackslash gon}, \texttt{\textbackslash set}, \texttt{\textbackslash torz}), with extra wrapper packages (bm, bbm, dsfont, xspace, cancel, mathtools, cleveref). The delivered numbering is the unit's own. Assets: no retained span contains a figure environment, a graphics-inclusion command, a PDF-page inclusion, a table or a TikZ picture; no image file is copied into this package, the package declares no asset dependency and no image file is redistributed with it. Licence: version arXiv:2412.16016v2 of this article is distributed under the Creative Commons Attribution 4.0 International licence, as recorded in the arXiv licence metadata for this exact version and shown on the abstract page https://arxiv.org/abs/2412.16016v2. A full rights notice is delivered at licenses/arxiv-2412.16016-CC-BY-4.0.txt. Characters: every retained excerpt is pure ASCII, so the delivered engine, XeLaTeX with its default Latin Modern text font, reports no missing character.
\begin{proposition}\label{prop:hecke-kernel}
Let $n$ be an integer, $q \nmid n$ a prime and let $H$ be a subgroup
  of~$(\Z/n\Z)^\times$ containing $-1$. Let $a=\sum_{i=1}^k a_i\diamondop{d_i} \in \Z[(\Z/n\Z)^\times/H]$ be a linear combination of diamond operators. We consider $t :=T_q - a$ 
as a correspondence on $X_1(N)$, inducing an endomorphism of the divisor group of $X_1(N)$
over $\C$. Then the kernel of $t$ is contained in the subgroup of divisors supported in cusps.
\end{proposition}

\begin{lemma}\label{lem:special_cusp}
Let $m=1$ or $2$, $n$ an integer such that $m|n$ and $q$ a prime that does not divide $2n$. Let $C_0 \in X_1(m,n)(\mathbb \Z[1/n])$ be the cusp whose moduli interpretation is the  N\'eron $n$-gon $E_0 := \mathbb G_m \times \torz{n}$ together with the points $P_0:=((-1)^{m-1},0)$ of order $m$ and $Q_0 :=(1,1)$ of order $n$. Then $A_q(C_0)$ is $0$.
\end{lemma}

\begin{proposition}[(Derickx)] \label{prop:hecke_gon}
  Let $d \ge 1$, let $n$ be an integer, $q \nmid n$ a prime and let $H$ be a subgroup
  of~$(\Z/n\Z)^\times$ containing $-1$. We assume that there is an $a \in (\Z/n\Z)^\times/ H$ such that $T := (\diamondop{a}-1)(J_H(\Q))$ is finite.
  When $q=2$ we either assume that $\#A$ is odd or, more specifically, that the $2$-primary part of~$T$ is killed by $A_2$.
  We then set
  \[ k_q = \begin{cases}
           q+1 & \text{if $J_H(\Q)$ is finite, otherwise} \\
           2q + 1 & \text{if $a \in \set{q, q^{-1}}$,} \\
           2(q + 1)& \text{if $a \notin \set{q, q^{-1}}$.}
         \end{cases}
  \]
  Then $k_qd < \gon_{\Q}(X_H)$ implies that any point
  in~$X_H^{(d)}(\Q)$ without cusps in its support is a sum of orbits
  under~$\diamondop{a}$. In particular, the inequality above holds when
  \begin{equation}\label{DS-inequality}
      d < \frac{325}{2^{15}} \, \frac{[\SL_2(\mathbb Z) :\Gamma_1(n)]}{k_q \cdot \#H}.
  \end{equation} 
\end{proposition}

\begin{proof}
Let $c \in X_1(N)(\Q)$ be a rational cusp such that $A_q(c) = 0$; such a cusp exists by \Cref{lem:special_cusp}.  Let $c'$ be the image of $c$ in $X_H(\Q)$. Define 

\[ B_{a} :=\begin{cases}
           1 & \text{if $J_H(\Q)$ is finite, } \\
           (\diamondop{a}-1) & \text{otherwise.}
         \end{cases}
         \]
Let $D \in X_H^{(d)}(\Q)$ be an effective divisor of degree $d$ without cusps in its support, and consider the linear equivalence class $[D-dc']$ as a rational point in $J_H(N)(\Q)$. Then $A_qB_{a}([D-dc'])$ is $0$. Indeed, if $J_H(\Q)$ is finite then $A_qB_{a}([D-dc'])=0$ because $A_q$ kills all the torsion. If $J_H(\Q)$ is not finite, then $(\diamondop{a}-1)([D-dc'])$ is torsion by assumption and hence also killed by $A_q$.

Now since $A_qc' =0$ it follows that $A_qB_{a}(D-dc')=A_qB_{a}D$.
Now if $[A_qB_{a}D]=0$ there are two possibilities, either $A_qB_{a}D=0$ or there is a non-constant function $f$ such that $A_qB_{a}D = \ddiv f$.

First consider the case $A_qB_{a}D=0$. Since $D$ and hence $B_{a}D$ does not contain any cusps in its support, one can apply \Cref{prop:hecke-kernel} to deduce $B_{a}D=0$. If $J_H(\Q)$ is finite, this implies $D=0$, so in this case there does not exist an effective divisor $D \in X_H^{(d)}(\Q)$ of degree $d$ without cusps in its support, and there is nothing to prove. If $J_H(\Q)$ is not finite, then $B_{a}D=0$ implies $\diamondop a D =D$ forcing $D$ to be a sum of orbits under $\diamondop a$.

What remains to show is that the second case, $A_qB_{a}D = \ddiv f$ for some non-constant function $f$, cannot happen.

By separating the positive and the negative terms one can rewrite $A_qB_{a}D$ as \begin{align}A_qB_{a}D=\begin{cases}
           T_qD - (q\diamondop{q}+1)D  & \text{if $J_H(\Q)$ is finite, } \\
           (T_q\diamondop{a}+q\diamondop{q}+1)D - (T_q +\diamondop{a}(q\diamondop{q}+1))D & \text{otherwise}. 
         \end{cases}\label{eq:linear-equivalence}\end{align}
First consider the case when $J_H(\Q)$ is finite. Let $k_q:=q+1$. Note that $T_qD$  and $(q\diamondop{q}+1)D$ are both effective divisors of degree $k_qd$, so $\deg f \leq k_qd$ which contradicts $k_qd < \gon_\Q(X_H).$

Now consider the case when $J_H(\Q)$ is not finite. Let $k_q:=2(q+1)$.
If $a \notin \set{q,q^{-1}}$ then $(T_q\diamondop{a}+q\diamondop{q}+1)D$ and $(T_q +\diamondop{a}(q\diamondop{q}+1))D$ are effective divisors of degree $k_qdd$ and $\deg f \leq k_qd$ again contradicting $k_qd < \gon_\Q(X_H).$
In the cases $a=q$ and $a=\diamondop{q^{-1}}$ the right-hand side of \eqref{eq:linear-equivalence} simplifies to $(T_q\diamondop{q}+(q-1)\diamondop{q}+1)D - (T_q +q\diamondop{q^2})D$ respectively $ (T_q\diamondop{a}+q\diamondop{q})D - (T_q +(q-1)\diamondop{q}+\diamondop{q^{-1}})D.$ This simplification decreases the degree of the positive and negative part by 1 and hence there is again a contradiction with $k_qd < \gon_\Q(X_H)$ as before.
         
\end{proof}

\end{document}
